\section{Design a Rate Limiter}
\label{sec:case-rate-limiter}

\problemstatement{Design a rate limiter that caps how many requests a
client may make in a time window -- e.g.\ 100 requests per minute -- and
lets the algorithm (token bucket, sliding window, fixed window) be swapped.
Return allow/deny per request.}

\begin{patternbox}[Clarify Before Designing]
Ask: Per-user, per-IP, or global limits? Which algorithm -- fixed window,
sliding window log/counter, token bucket, leaky bucket? Distributed across
servers or single-node? For this design assume: per-client keys, a
swappable algorithm, and single-node with a note on distributed
extensions.
\end{patternbox}

\subsection*{Identify the entities}

A \texttt{RateLimiter} facade keyed by client, and a
\texttt{RateLimitStrategy} interface with implementations
(\texttt{TokenBucket}, \texttt{FixedWindow}, \texttt{SlidingWindow}). Each
strategy stores per-client state and answers \texttt{allow()}. The varying
algorithm is, again, a Strategy.

\begin{ideabox}[Token Bucket in a Nutshell]
A token bucket holds up to \texttt{capacity} tokens, refilled at a steady
\texttt{rate} per second. Each request tries to take one token: succeed and
proceed, or fail and get rejected. It naturally allows short bursts (up to
capacity) while bounding the long-run average -- which is why it is the
most common choice. Refill lazily by computing elapsed time on each call
rather than running a timer.
\end{ideabox}

\subsection*{Class design (C++ skeleton)}

\begin{lstlisting}
#include <bits/stdc++.h>
using namespace std;

struct RateLimitStrategy {
    virtual bool allow(const string& client) = 0;
    virtual ~RateLimitStrategy() = default;
};

class TokenBucket : public RateLimitStrategy {
    double capacity, refillPerSec;
    struct Bucket { double tokens; double lastRefill; };
    unordered_map<string, Bucket> buckets;

    double now() {
        return chrono::duration<double>(
            chrono::steady_clock::now().time_since_epoch()).count();
    }
public:
    TokenBucket(double cap, double rate) : capacity(cap), refillPerSec(rate) {}
    bool allow(const string& client) override {
        double t = now();
        auto& b = buckets.try_emplace(client, Bucket{capacity, t}).first->second;
        b.tokens = min(capacity, b.tokens + (t - b.lastRefill) * refillPerSec);
        b.lastRefill = t;
        if (b.tokens >= 1.0) { b.tokens -= 1.0; return true; }   // consume
        return false;                                            // throttled
    }
};

class RateLimiter {                             // facade over the strategy
    unique_ptr<RateLimitStrategy> strategy;
public:
    RateLimiter(unique_ptr<RateLimitStrategy> s) : strategy(move(s)) {}
    bool allow(const string& client) { return strategy->allow(client); }
};

int main() {
    RateLimiter limiter(make_unique<TokenBucket>(3, 1));  // burst 3, 1/s
    for (int i = 0; i < 5; i++)
        cout << "req " << i << ": " << (limiter.allow("user1") ? "OK" : "DENY") << "\n";
    return 0;
}
\end{lstlisting}

\begin{takeawaybox}
Rate limiting is a \textbf{Strategy} problem: token bucket, fixed window,
and sliding window are interchangeable behind one \texttt{allow()}
interface. Know token bucket's burst-vs-average trade-off cold, and
compute refill lazily. The senior-level extension is \emph{distributed}
limiting -- shared counters in Redis with atomic increments and expiry.
\end{takeawaybox}
